\chapter{Diagramas UML de casos de uso}
\label{apendice-casos-uso}

Este apêndice apresenta os diagramas UML correspondentes aos principais casos de uso do sistema Biribá. As funcionalidades foram agrupadas por área para evidenciar as associações dos atores agricultor, coordenador e administrador com o sistema, sem omitir nenhum dos casos de uso identificados no levantamento de requisitos.

\tikzset{
    uml system/.style={draw, rounded corners=2pt, fill=gray!4},
    uml use case/.style={
        draw,
        ellipse,
        fill=white,
        align=center,
        font=\small,
        minimum width=4.4cm,
        minimum height=0.95cm
    },
    uml association/.style={draw, semithick},
    pics/uml actor/.style args={#1}{
        code={
            \draw (0,0.65) circle (0.17);
            \draw (0,0.48) -- (0,-0.12);
            \draw (-0.32,0.27) -- (0.32,0.27);
            \draw (0,-0.12) -- (-0.27,-0.58);
            \draw (0,-0.12) -- (0.27,-0.58);
            \node[font=\small, align=center, anchor=north] at (0,-0.72) {#1};
            \coordinate (-center) at (0,0.15);
        }
    }
}

\section{Acesso ao sistema}

A Figura~\ref{fig:uc-acesso} representa a solicitação de cadastro pelo agricultor e a aprovação ou rejeição desse acesso pelo coordenador ou administrador.

\begin{figure}[H]
    \centering
    \begin{tikzpicture}
        \node[uml system, minimum width=8.4cm, minimum height=4.2cm] (sistema) at (7.5,0) {};
        \node[anchor=north west, font=\small\bfseries] at ([xshift=0.2cm,yshift=-0.15cm]sistema.north west) {Sistema Biribá};

        \node[uml use case] (solicitar) at (7.5,1) {Solicitar cadastro};
        \node[uml use case] (aprovar) at (7.5,-1) {Aprovar usuário};

        \pic (agricultor) at (1.4,1) {uml actor={Agricultor}};
        \pic (coordenador) at (1.4,-1) {uml actor={Coordenador}};
        \pic (administrador) at (13.6,-1) {uml actor={Administrador}};

        \draw[uml association] (agricultor-center) -- (solicitar.west);
        \draw[uml association] (coordenador-center) -- (aprovar.west);
        \draw[uml association] (administrador-center) -- (aprovar.east);
    \end{tikzpicture}
    \caption{Diagrama UML dos casos de uso de acesso ao sistema}
    \label{fig:uc-acesso}
\end{figure}

\section{Configuração do edital}

A Figura~\ref{fig:uc-configuracao} apresenta os casos de uso associados à preparação da chamada pública. O coordenador e o administrador podem gerenciar famílias, criar editais e criar ciclos de entrega, enquanto o gerenciamento global de produtos e destinos é atribuído ao administrador.

\begin{figure}[H]
    \centering
    \begin{tikzpicture}
        \node[uml system, minimum width=8.6cm, minimum height=7.2cm] (sistema) at (7.5,0) {};
        \node[anchor=north west, font=\small\bfseries] at ([xshift=0.2cm,yshift=-0.15cm]sistema.north west) {Sistema Biribá};

        \node[uml use case] (familias) at (7.5,2.35) {Gerenciar famílias};
        \node[uml use case] (edital) at (7.5,1.18) {Criar edital};
        \node[uml use case] (produtos) at (7.5,0) {Gerenciar produtos};
        \node[uml use case] (destinos) at (7.5,-1.18) {Gerenciar destinos};
        \node[uml use case] (ciclos) at (7.5,-2.35) {Criar ciclos de entrega};

        \pic (coordenador) at (1.3,0.7) {uml actor={Coordenador}};
        \pic (administrador) at (13.7,-0.1) {uml actor={Administrador}};

        \draw[uml association] (coordenador-center) -- (familias.west);
        \draw[uml association] (coordenador-center) -- (edital.west);
        \draw[uml association] (coordenador-center) -- (ciclos.west);

        \draw[uml association] (administrador-center) -- (familias.east);
        \draw[uml association] (administrador-center) -- (edital.east);
        \draw[uml association] (administrador-center) -- (produtos.east);
        \draw[uml association] (administrador-center) -- (destinos.east);
        \draw[uml association] (administrador-center) -- (ciclos.east);
    \end{tikzpicture}
    \caption{Diagrama UML dos casos de uso de configuração do edital}
    \label{fig:uc-configuracao}
\end{figure}

\section{Execução e acompanhamento}

A Figura~\ref{fig:uc-execucao} representa as funcionalidades utilizadas durante a execução do edital. O coordenador e o administrador registram entregas, acompanham os indicadores e geram relatórios, enquanto o agricultor consulta os dados referentes à sua participação.

\begin{figure}[H]
    \centering
    \begin{tikzpicture}
        \node[uml system, minimum width=8.6cm, minimum height=6.3cm] (sistema) at (7.5,0) {};
        \node[anchor=north west, font=\small\bfseries] at ([xshift=0.2cm,yshift=-0.15cm]sistema.north west) {Sistema Biribá};

        \node[uml use case] (entrega) at (7.5,1.8) {Registrar entrega};
        \node[uml use case] (dashboard) at (7.5,0.6) {Acompanhar dashboard};
        \node[uml use case] (relatorios) at (7.5,-0.6) {Gerar relatórios};
        \node[uml use case] (consulta) at (7.5,-1.8) {Consultar informações\\da participação};

        \pic (coordenador) at (1.3,0.6) {uml actor={Coordenador}};
        \pic (administrador) at (13.7,0.6) {uml actor={Administrador}};
        \pic (agricultor) at (1.3,-1.8) {uml actor={Agricultor}};

        \draw[uml association] (coordenador-center) -- (entrega.west);
        \draw[uml association] (coordenador-center) -- (dashboard.west);
        \draw[uml association] (coordenador-center) -- (relatorios.west);

        \draw[uml association] (administrador-center) -- (entrega.east);
        \draw[uml association] (administrador-center) -- (dashboard.east);
        \draw[uml association] (administrador-center) -- (relatorios.east);

        \draw[uml association] (agricultor-center) -- (consulta.west);
    \end{tikzpicture}
    \caption{Diagrama UML dos casos de uso de execução e acompanhamento}
    \label{fig:uc-execucao}
\end{figure}

\section{Administração da plataforma}

A Figura~\ref{fig:uc-administracao} apresenta o caso de uso de manutenção da plataforma, reservado ao administrador e relacionado aos ajustes, às correções e à evolução do sistema.

\begin{figure}[H]
    \centering
    \begin{tikzpicture}
        \node[uml system, minimum width=8.4cm, minimum height=3cm] (sistema) at (7.5,0) {};
        \node[anchor=north west, font=\small\bfseries] at ([xshift=0.2cm,yshift=-0.15cm]sistema.north west) {Sistema Biribá};

        \node[uml use case] (manter) at (7.5,0) {Manter o sistema};
        \pic (administrador) at (1.4,0) {uml actor={Administrador}};

        \draw[uml association] (administrador-center) -- (manter.west);
    \end{tikzpicture}
    \caption{Diagrama UML do caso de uso de administração da plataforma}
    \label{fig:uc-administracao}
\end{figure}
